\textbf{Auto-Regressive Language Models (ARLMs)} Large language models \cite{devlin2019bertpretrainingdeepbidirectional} are
    given a vocabulary of $\sum$ and text sequences \(\omega = (w_1,\dots,w_T) \in \Sigma^T\). In an auto-regressive model, $\theta$
    denotes the parameters that define a conditional probability distribution over the next token given its predecessors,
  $p_\theta(w_t
    | w_{1:t-1})$. The training objective is to choose \(\theta\) so that \(p_\theta\) approximates the data distribution
    \(p_{\text{data}}\) by minimizing the expected negative log-likelihood.

  \textbf{Diffusion Language Models (DLMs)} For diffusion language models, the goal is to learn the conditional
  distributions that will invert a corruption process. The model aims to transform a masked, noisy output sequence to samples from \(p_{\text{data}}\). The goal is to minimize a cross entropy loss only on the masked tokens. At each refinement step, the model unmasks its most confident predictions while re-masking the tokens with the lowest confidence for subsequent processing. These models have been implemented in several prominent systems, such as LLaDA, and similar diffusion-based approaches have been explored in models like Gemini Diffusion.
  \newline
  \textbf{LLM Watermark}
  A watermark is a hidden statistical signal embedded in machine-generated text to determine if it is AI generated. A watermark changes the model’s word choices during generation, guiding it towards a particular pattern that would be unlikely to appear randomly. Changes made by the watermark are designed to be imperceptible to a human reader, preserving text quality and coherence. However, a
  detection algorithm with access to a secret key can analyze the text and detect a pattern, verifying that the text was generated by a
  watermarked model.

   A language model watermark can be formally defined as a set of three components: an embedding process \(E\), a sampling procedure
  \(G\), and a
    detection function \(D\). These components operate on the underlying token distributions of a given language model.

    Let \(p_\theta(w_t | c)\) be the original conditional probability distribution over the vocabulary \(\Sigma\) that the language
  model \(\theta\)
    assigns to the token \(w_t\) given a relevant context \(c\).

      1. Embedding Process (\(E\)): The embedding process transforms the original distribution \(p_\theta\) into a watermarked
         distribution \(p'_\theta\).
        This transformation is conditioned on a message or secret key \(k\).
        \[ E(p_\theta(w_t | c), k) \rightarrow p'_\theta(w_t | c, k) \]
        To preserve the quality of the generated text, the watermarked distribution must remain close to the original. This is
  formally constrained by
    requiring that a statistical distance, such as the Kullback-Leibler (KL) divergence, between the two distributions is bounded by a
  small value
    \(\epsilon\):
        \[ D_{KL}(p'_\theta || p_\theta) \leq \epsilon \]

      2. Sampling Procedure (\(G\)): The generation or sampling procedure produces a watermarked text \(\omega = (w_1, \dots, w_T)\)
         by drawing tokens
        from the watermarked distribution \(p'_\theta\). For autoregressive models, this involves sequential sampling: \(w_t \sim
  p'_\theta(w_t |
        w_{1:t-1}, k)\). For Diffusion Language Models (DLMs), this involves an iterative refinement process where the model
  repeatedly samples and
        updates tokens across the sequence, guided by \(p'_\theta\) to invert a corruption process. In both cases, the context \(c\)
  in \(p'_\theta(w_t |
        c, k)\) represents the relevant information at each decision point.

      3. Detection Function (\(D\)): The detection function determines if a given text \(\omega\) contains the watermark. This is a
         hypothesis testing
        problem with two hypotheses:
          Null Hypothesis (\(H_0\)): The text was generated by the original model, \(\omega \sim p_\theta\).
          Alternative Hypothesis (\(H_1\)): The text was generated by the watermarked model, \(\omega \sim p'_\theta\).

        The detector computes a test statistic \(S(\omega, k)\) and compares it to a threshold \(\tau\). The text is classified as
  watermarked if
    \(S(\omega, k) > \tau\). The choice of \(\tau\) determines the trade-off between the False Positive Rate (Type I error,
  \(\alpha\)) and the True
    Positive Rate (power, \(1-\beta\)).

  \textbf{Properties of an Effective Watermark}
  An effective watermarking scheme is expected to have several key properties:
  \begin{itemize}
      \item \textbf{Text Quality Preservation} The watermark should not clearly degrade the quality, fluency, or correctness of
  generated text. The KL-divergence constraint on \(p'_\theta\) is a formal measure of this.
      \item \textbf{Robustness:} The statistical signal must be robust to common text alterations. A user attempting to remove the
  watermark might paraphrase sentences, insert or delete words, or use another model to revise the text. The detection function \(D\)
  should still identify the watermark with after such modifications.
      \item \textbf{Security:} It should be computationally infeasible to detect the watermark, forge it in human-written
  text, or remove it from generated text without access to the secret key \(k\).
      \item \textbf{Reliability:} The watermark must be reliably detectable, achieving a high True Positive Rate while maintaining a low False Positive Rate to avoid falsely accusing human authors.
  \end{itemize}

  \textbf{Prior Work in LLM Watermarking}
  The most prominent watermarking scheme for ARLMs \cite{kirchenbauer2023} partitions the vocabulary into a "green list" and a "red list" at each time step. The partition is determined pseudorandomly using a secret key and the hash of the preceding token sequence. The embedding process \(E\) slightly increases the probabilities of tokens on the green list before sampling. The detection function \(D\) involves counting the number of green-listed tokens in a given text to determine whether it statistically significant. While effective, this reliance on a fixed context window makes it incompatible with the non-sequential generation of DLMs. Early work in watermarking DLMs[citation idk how to do pls help] has attempted to adapt this method by computing context hashes in expectation over the masked tokens.